\chapter{对齐与强化学习：RLHF、PPO 与 NLP 中的 RL}

\begin{introduction}
\item {\heiti 本章核心问题}：除了让模型“答对”，还应该怎样让它“按希望的方式回答”，以及 RLHF 在这个过程中到底解决什么问题？
\item {\heiti 主案例推进到哪一步}：把企业知识助手从“会完成任务”推进到“更符合偏好与风险约束”。
\item {\heiti 读完后能做什么}：知道什么时候 RLHF 值得做，什么时候它只是掩盖基础问题的错误方向；也能理解 DPO、RLAIF 等替代路线为何在很多团队里更现实。
\end{introduction}

\section{学习目标}
\begin{itemize}
  \item 理解为什么仅靠 SFT 仍不足以完整控制模型行为。
  \item 掌握强化学习在大模型中的最小概念集，知道奖励、策略、轨迹和策略更新分别在讲什么。
  \item 掌握 RLHF 的基本流程，以及 PPO 和 KL 约束在其中的角色。
  \item 能设计基础的偏好数据方案，尤其能为企业知识助手写出可执行的偏好样例。
  \item 能识别奖励模型失败模式、奖励黑客和“现在不该上 RLHF”的工程信号。
\end{itemize}

\section{问题背景}
SFT 可以让模型学会“像训练答案那样输出”，但真实系统的目标往往不止是答对。团队还会关心更复杂的行为：回答是否更安全、更简洁、更保守，是否优先给出依据，是否在信息不足时主动澄清，是否能在高风险场景下把决策升级给人工。上述要求并不总能用一份唯一标准答案表达出来，因此出现了对齐训练的需求。

从工程角度看，对齐问题的本质不是“模型会不会做题”，而是“模型在多种可行答案里，是否更偏向我们接受的那一种”。企业知识助手尤其如此。很多时候，不同回答在事实层面都不算错，但在保守性、引用规范、风险提示、权限边界和行动建议上差异很大。只靠标准 SFT 样本，很难把这种偏好完整编码进去。

上一章我们已经把评测门禁搭起来了，现在才谈得上对齐。因为如果团队连“哪些回答更好、哪些退化不能接受”都看不清，任何对齐训练都只是在放大主观感觉。再进一步说，很多失败案例并不是模型无知，而是模型在“会回答”之后仍然作出了错误风格选择。例如它可能知道制度条文，却没有先引用依据；也可能知道自己证据不足，却依旧给出过度肯定的建议。此时如果继续追加普通 SFT 数据，往往只会堆出更多模板，而不一定能让模型真正学会在多个候选回答之间作出更符合业务目标的选择。

\section{核心概念}
\subsection{为什么需要对齐}
SFT 的优化目标是最大化标准答案的似然，但真实偏好往往是相对性的。人类更容易判断“两段回答哪一段更好”，而不是为每个问题写出唯一正确答案。对齐训练就是把这种相对偏好引入模型优化过程。

可以把模型优化目标分成三层来看：
\begin{itemize}
  \item \textbf{能力层}：模型会不会做这件事。
  \item \textbf{行为层}：模型做这件事时，是否按指定风格和边界输出。
  \item \textbf{策略层}：在多种可行行为中，模型优先选择哪一种。
\end{itemize}

SFT 主要解决前两层，而 RLHF 更偏向解决第三层。也正因如此，它通常不应替代基础能力建设，而应建立在基础能力已经足够可用的前提上。一个检索错误、知识过期、工具调用失败的系统，并不会因为加了 RLHF 就变成一个可靠系统；它只会更稳定地表现出原本的问题。

对齐之所以必要，还因为真实业务偏好经常互相拉扯。团队既希望模型回答快，又希望它谨慎；既希望它简洁，又希望它可追溯；既希望它帮用户推进任务，又不能越权替用户作决定。对齐训练不是让模型学会一条绝对规则，而是让它在冲突目标之间更稳定地落在团队能接受的位置上。

\subsection{RL 的最小概念集}
对工程读者来说，不需要把强化学习学成完整课程，但至少要建立四个最小概念：
\begin{itemize}
  \item \textbf{策略}：模型当前的行为规则，也就是“给定输入时倾向输出什么”。
  \item \textbf{奖励}：系统对行为结果的偏好信号。
  \item \textbf{轨迹}：一次输入到输出的完整过程。
  \item \textbf{策略更新}：根据奖励调整模型，让高奖励行为更容易再次出现。
\end{itemize}

在大语言模型里，这些概念并不是拿来玩游戏，而是用来回答一个更实务的问题：\textbf{模型已经会答了，怎样让它更稳定地答成我们想要的样子。}

如果继续压缩成“RL 最小集”，可以记成一句话：模型先生成候选回答，再由奖励信号评价这些回答，最后通过策略更新让“更受偏好”的回答概率上升。这里最难的从来不是“更新”本身，而是“奖励到底代表什么”。奖励定义模糊，后面的学习过程就会稳定地把模糊放大。

\subsection{RLHF 的基本流程}
典型 RLHF 流程包括三步：
\begin{enumerate}
  \item 先做 SFT，得到能完成任务的基础策略。
  \item 收集偏好数据，训练奖励模型，学习人类对不同回答的偏好。
  \item 用强化学习方法优化策略，使模型倾向于产生高奖励回答。
\end{enumerate}

这里最关键的转折点，在于训练信号从“标准答案是什么”变成了“这两个答案哪个更好”。这让模型开始学习偏好，而不仅仅是学习模板。

很多教材会把 RLHF 讲成一条很顺的流水线，但实际项目常常不是一次走完，而是循环迭代：先从线上失败案例抽取偏好样本，训练一个初版奖励模型，再做一轮小步策略更新，然后回到评测集与人工抽检，确认收益来自“更好”，而不是来自“更会讨好打分器”。因此 RLHF 更像一个持续校准系统，而不是一把单次升级的锤子。

\subsection{高质量 SFT 仍是对齐起点：LIMA 的启示不是“样本越少越神”}
旧稿里提到过 LIMA，这个案例在教材里最值得保留的，不是“1000 条样本也能做出不错效果”这句口号，而是它提醒读者：\textbf{很多对齐收益其实先来自高质量交互数据，而不是来自后面的 RL 名字。}

如果基础 SFT 已经把格式、引用顺序、拒答语气和任务边界训得很稳，那么后续对齐只需要处理“多个都还可以的答案里，哪一个更符合偏好”。反过来，如果基础 SFT 还没把最基本的交互协议训稳，直接进入 RLHF 往往会让团队在更昂贵的链路里追一个本来更便宜就能解决的问题。

LIMA 给工程团队的真正启示可以压缩成两句：第一，\textbf{预训练解决知识与能力底座，SFT 先解决交互样式和基础行为。} 第二，\textbf{只有当偏好问题真的剩下来时，RLHF、DPO 一类方法才开始值钱。} 这也是为什么成熟团队常把“是否已经有一版高质量 SFT 基线”当成进入对齐训练前的硬门槛。

\subsection{偏好数据与奖励模型}
奖励模型并不是凭空出现的，它来自偏好数据。偏好数据最常见的形式，不是长篇标注说明，而是成对比较：给定同一个问题，让标注者选择答案 A 和答案 B 谁更好。与其要求标注者写“完美答案”，很多场景下更现实的做法是要求他们判断“哪一个更符合团队偏好”。

\begin{table}[H]
\centering
\small
\begin{tabularx}{\textwidth}{>{\raggedright\arraybackslash}p{3cm}>{\raggedright\arraybackslash}p{3.6cm}>{\raggedright\arraybackslash}X}
\toprule
\textbf{数据形式} & \textbf{更容易表达什么} & \textbf{主要风险} \\
\midrule
标准 SFT 答案 & 任务格式、固定写法、基础行为 & 容易把偏好问题压成单一模板 \\
成对偏好比较 & 哪个答案更保守、更有依据、更符合风格 & 标注标准不清时容易漂移 \\
线上反馈回流 & 真实使用中的好坏判断 & 噪声大，且容易受界面、上下文和用户情绪干扰 \\
\bottomrule
\end{tabularx}
\caption{SFT 数据、偏好数据和线上反馈的职责不同}
\label{tab:preference-data-types}
\end{table}

对企业知识助手来说，偏好数据特别适合表达下面这类差异：
\begin{itemize}
  \item 同样都引用了材料，但哪个答案更保守。
  \item 同样都答对了，哪个答案更先给结论再给依据。
  \item 同样都不完整，哪个答案更愿意明确说“材料不足”。
  \item 同样都给出建议，哪个答案更遵守权限边界，不替用户做越权判断。
\end{itemize}

偏好数据设计的第一原则不是“尽量多”，而是“偏好维度明确”。如果同一批样本里同时混入“更简洁”“更详细”“更大胆推进”“更保守拒答”四种互相冲突的目标，奖励模型就会学到一个平均化、不可解释的标准。更好的做法是先写清偏好维度，再决定哪些维度可以混合标注，哪些需要单独分层。

以企业知识助手为例，可以把偏好维度先拆成以下几类：\textbf{证据优先、保守性、可执行性、权限边界、表达风格}。下面给出几个更具体的偏好样例。

\begin{table}[H]
\centering
\small
\begin{tabularx}{\textwidth}{>{\raggedright\arraybackslash}p{2.7cm}>{\raggedright\arraybackslash}X>{\raggedright\arraybackslash}X}
\toprule
\textbf{偏好维度} & \textbf{优选回答特征} & \textbf{不优选回答特征} \\
\midrule
证据优先 & 先给出制度条款或知识库出处，再给结论与动作建议 & 直接下判断，不说明依据来自哪里 \\
保守性 & 证据不足时明确说明不确定性，并给出补充信息需求 & 材料不全时仍然给出肯定结论 \\
权限边界 & 对审批、财务、人事等高风险事项，建议人工复核或升级流程 & 直接替用户决定“可以报销”“可以解雇”“可以上线” \\
可执行性 & 建议能落到下一步动作，如“需补齐字段、补传附件、联系审批人” & 只有原则性空话，没有下一步动作 \\
表达风格 & 语言简洁、条理清楚、便于工单系统记录 & 话很多但结构混乱，或过度套话、显得很像免责声明模板 \\
\bottomrule
\end{tabularx}
\caption{企业知识助手中常见的偏好维度}
\label{tab:enterprise-preference-dimensions}
\end{table}

再看一个成对偏好例子。用户提问：“员工离职当天是否还能访问客户资料库？”若知识库中只有“离职前需完成资产交接”和“客户资料访问需按岗位授权”两条制度，但没有明确写“离职当天”的例外处理，那么更好的回答通常不是武断地说“不能”或“可以”，而是：
\begin{quote}
现有制度能支持的结论是：客户资料访问受岗位授权控制，但你提供的材料没有明确说明离职当天权限何时回收。更稳妥的处理是先联系 IT 权限管理员或查离职流程单中的权限回收节点；若你要执行操作，请不要仅依据当前回答直接放行。
\end{quote}

之所以偏好这类回答，不是因为它更长，而是因为它同时满足三件事：说明依据、承认缺口、给出下一步动作。反过来，“离职当天一律不能访问客户资料库，请立即封禁账号”虽然看起来果断，但如果制度里并没有这一条，它就是一个高风险偏好错误。

奖励模型训练时还要处理一个常见现实：\textbf{标注共识远比标注速度重要}。如果两个标注员对同类问题的胜负判断经常相反，说明偏好标准还没有定清楚。这时继续扩样本量意义不大，更应该回到标注手册，明确优先级，例如“证据优先高于措辞礼貌”“权限边界高于任务完成度”“高风险场景宁可保守，不奖励过度推进”。

\subsection{Bradley--Terry：奖励模型怎样把“哪个更好”变成可学习分数}
旧稿里给过奖励模型训练的一个关键公式，这里值得保留，因为它恰好能把“人类偏好”转成工程上可训练的对象。假设对同一个问题 \(x\)，回答 \(y_w\) 比回答 \(y_l\) 更受偏好，Bradley--Terry 风格的建模方式会把这种胜负写成：
\[
P(y_w \succ y_l \mid x)=\sigma\left(r_{\phi}(x,y_w)-r_{\phi}(x,y_l)\right),
\]
其中 \(r_{\phi}(x,y)\) 是奖励模型对回答的打分，\(\sigma(\cdot)\) 是 sigmoid 函数。

这个式子的工程直觉其实很简单：\textbf{奖励模型不用先学“绝对满分是什么”，它先学“两个答案比起来谁更好”。} 只要它能在足够多的比较对上把更优回答分得更高，后面的策略优化就有了方向。

这也是为什么偏好数据质量如此关键。如果比较对本身就带着歧义，或者标注标准前后不一，那么奖励模型学到的差值就不是团队真实偏好，而是标注噪声。到后面不管你接 PPO 还是 DPO，都会把这种噪声放大。因此，Bradley--Terry 公式最重要的教学意义不是数学优雅，而是提醒读者：\textbf{奖励建模首先是比较协议建模。}

\subsection{奖励模型不是黑箱评分器：先把标签协议和模型角色分清}
旧稿里还讲过奖励模型训练和一致性选择。教材里最值得保留的一层，是把奖励模型当成一个\textbf{有评分协议的近似裁判}，而不是一个“会自动懂业务”的黑箱。

训练奖励模型之前，团队最好先写清三件事：
\begin{itemize}
  \item \textbf{比较单位是什么}：是完整回答的整体优劣，还是证据性、保守性、可执行性中的某一维。
  \item \textbf{胜负规则是什么}：是否允许平局，是否允许“都不好但 A 稍好”，是否要对高风险场景单独加权。
  \item \textbf{输入材料包含什么}：只看问题和回答，还是连同检索证据、系统提示和风险等级一起给奖励模型。
\end{itemize}

如果这三件事不清，奖励模型很容易学成“语言风格评分器”，而不是“业务偏好评分器”。例如，只给它看问题和回答，不给它看证据片段，它就很难真正判断“引用是否支持结论”；只给它一个总分胜负，不区分高风险和低风险场景，它就更容易把“说得顺”误当成“说得对”。

\subsection{偏好样本也有格式：比较对之外，还要带上风险与证据}
旧稿在奖励模型数据格式上讲得比当前重写版更细，这部分不该丢。教材里至少要让读者知道：\textbf{偏好样本不是“题目 + A/B 两个回答 + 胜者”就结束了。} 如果要让奖励模型学到的是业务偏好，而不是表面文风，样本字段本身就要能承载业务约束。

\begin{table}[H]
\centering
\small
\begin{tabularx}{\textwidth}{>{\raggedright\arraybackslash}p{2.8cm}>{\raggedright\arraybackslash}p{3.2cm}>{\raggedright\arraybackslash}X}
\toprule
\textbf{偏好样本字段} & \textbf{为什么要有} & \textbf{缺了会怎样} \\
\midrule
用户问题 / 任务指令 & 给奖励判断一个明确任务边界 & 无法区分“答非所问”与“答得保守” \\
候选回答 A/B & 形成可比较对象 & 只能做单点评分，难以稳定学习相对偏好 \\
证据片段或上下文 & 判断结论是否被材料支持 & 奖励模型更容易学成措辞评分器 \\
胜负标签 / 平局标签 & 明确比较结果 & 无法表示“两者都差”或“差异不大” \\
风险等级 / 任务类别 & 让高风险样本获得单独尺度 & 高风险和低风险问题会被混成同一种偏好 \\
标注理由摘要 & 便于质检和分歧分析 & 胜负相反时难以回溯是哪条标准冲突了 \\
\bottomrule
\end{tabularx}
\caption{偏好数据如果缺字段，奖励模型学到的往往不是团队真正想要的偏好}
\label{tab:preference-sample-schema}
\end{table}

这张表的意义不在于把数据结构做复杂，而在于提醒读者：偏好学习同样是一种数据工程。尤其在企业场景里，如果候选回答是否引用了制度、是否承认不确定性、是否触发人工升级，本来就属于不同维度的判断，那么数据字段里最好显式保留这些上下文，而不是指望奖励模型从一句“哪个好”里自动猜出来。

\subsection{奖励模型看什么、不看什么，要先把边界立住}
偏好数据一旦进入奖励模型，很多团队会不自觉地产生一个危险错觉：既然奖励模型最后要给回答打分，那是不是所有“系统希望它更好”的东西都可以交给奖励模型去学？答案是否定的。奖励模型擅长学习的是\textbf{相对偏好}，不擅长单独承担\textbf{事实校验、权限裁决和业务真值维护}。

\begin{table}[H]
\centering
\small
\begin{tabularx}{\textwidth}{>{\raggedright\arraybackslash}p{3cm}>{\raggedright\arraybackslash}p{3cm}>{\raggedright\arraybackslash}X}
\toprule
\textbf{信号类型} & \textbf{是否适合主要交给奖励模型} & \textbf{原因} \\
\midrule
回答是否更有依据、更保守、更清楚 & 适合 & 这是典型的相对偏好判断，人类或 AI 评审更容易比较 \\
制度事实是否为最新版本 & 不适合单独交给它 & 这更依赖知识库与版本治理，而不是回答表面风格 \\
权限是否真的允许执行动作 & 不适合单独交给它 & 权限边界应由鉴权与工作流系统硬性托底 \\
工具调用结果是否真实、字段是否齐全 & 不适合单独交给它 & 这应由工具侧校验和结构化检查负责 \\
不同风险级别下是否需要升级人工 & 适合部分学习，但必须配合规则系统 & 这是偏好与规则交叉区域，不能只靠模型自行体会 \\
\bottomrule
\end{tabularx}
\caption{奖励模型最擅长学习“比较标准”，不擅长替代业务真值系统}
\label{tab:reward-model-boundary}
\end{table}

这张表背后的工程判断非常重要。奖励模型当然可以帮助模型学会“高风险题更保守”，但它不应该成为你们权限系统的替身；它可以奖励“引用了证据且承认缺口”的回答，但不应该替代知识库版本控制去保证证据真伪。否则团队很容易把原本该由外部系统硬性保证的边界，错误地下放给一个统计学习模块。

\subsection{奖励模型也要校准：先看分歧样本，再看平均分}
奖励模型训练完成，并不意味着它已经“懂了团队偏好”。它更像一个刚被训练好的评分员，接下来还必须被校准。对工程团队来说，最有价值的校准样本通常不是那些胜负一眼就能看出来的简单题，而是\textbf{人和模型容易分歧、不同标注员也容易分歧的边界题}。这些题最能暴露奖励模型到底学到的是业务标准，还是表面语气。

一个实用的奖励模型校准流程可以很简单：先保留一小批金标比较对，不参与训练；训练后先看这批样本上的胜负一致性；再单独看高风险切片、证据不足切片和跨文档切片；最后把模型最自信但人类最不认同的样本拿出来做复盘。只有这三步都过得去，奖励模型才值得进入后续策略优化。

\begin{table}[H]
\centering
\small
\begin{tabularx}{\textwidth}{>{\raggedright\arraybackslash}p{2.8cm}>{\raggedright\arraybackslash}p{3cm}>{\raggedright\arraybackslash}X}
\toprule
\textbf{校准检查项} & \textbf{主要在看什么} & \textbf{一旦出问题通常说明什么} \\
\midrule
保留金标胜负一致性 & 奖励模型是否能复现已知偏好 & 基本比较协议或训练样本质量有问题 \\
高风险切片单独复核 & 模型是否把高风险题当成另一种尺度 & 风险等级没有真正进入奖励建模 \\
分歧样本复盘 & 人和模型最不一致的题集中在哪类规则 & rubric 优先级还没写清，或样本字段缺信息 \\
上线后坏例回放校验 & 新问法、新材料形态下是否快速漂移 & 奖励模型分布外脆弱，需补样本再校准 \\
\bottomrule
\end{tabularx}
\caption{奖励模型也需要像评测系统一样做持续校准}
\label{tab:reward-model-calibration}
\end{table}

教材里要特别强调一个反直觉点：\textbf{奖励模型的危险，不只是“打分不准”，而是“它会把不准稳定放大”。} 因为一旦后面接上 PPO 或 DPO，策略模型会主动向这种偏差靠拢。所以在对齐链路里，分歧样本往往比高分样本更值得团队反复看。

\subsection{RLHF 里的四个模型各自负责什么}
很多读者第一次接触 RLHF 时，最容易混淆的不是公式，而是到底有哪些模型在参与。旧稿里提到过 PPO 同时需要多个角色模型，这一点在新版里也应该说透。一个典型 RLHF/PPO 训练环里，常见有四个角色：
\begin{itemize}
  \item \textbf{策略模型（policy）}：当前正在被优化的回答者。
  \item \textbf{参考模型（reference）}：提供 KL 约束的参照，防止策略跑偏太远。
  \item \textbf{奖励模型（reward model）}：负责给生成结果打偏好分。
  \item \textbf{价值模型（value model）}：估计当前状态下未来奖励，用于更稳定地计算优势。
\end{itemize}

\begin{table}[H]
\centering
\small
\begin{tabularx}{\textwidth}{>{\raggedright\arraybackslash}p{2.8cm}>{\raggedright\arraybackslash}p{3cm}>{\raggedright\arraybackslash}X}
\toprule
\textbf{模型角色} & \textbf{主要职责} & \textbf{一旦出问题会怎样} \\
\midrule
策略模型 & 生成回答并被持续更新 & 回答行为本身直接漂移 \\
参考模型 & 提供“别偏太远”的基线 & KL 失去锚点，策略更容易冲偏 \\
奖励模型 & 对候选回答打偏好分 & 会把错误偏好稳定放大成训练目标 \\
价值模型 & 帮助 PPO 稳定估计优势 & 更新噪声变大，训练更容易抖动 \\
\bottomrule
\end{tabularx}
\caption{RLHF/PPO 里常见的四个模型角色}
\label{tab:rlhf-four-models}
\end{table}

这张表的工程含义在于：当 RLHF 训练效果不对时，不要只盯着策略模型。很多问题其实更早发生在奖励模型或价值估计上，只是最后由策略模型把问题显性化了。

\subsection{奖励模型与策略模型不必完全同构，但输入边界必须一致}
旧稿里提到过奖励模型与基础模型的一致性问题，这部分值得以更工程化的方式保留下来。教材里不需要把它讲成架构争论，但至少要让读者知道：\textbf{奖励模型可以不和策略模型完全同构，可它们看到的任务边界最好是一致的。}

所谓“一致”，不一定是参数规模一致、层数一致，而是下面几件事最好尽量统一：同一套 tokenizer 或至少同一套文本切分假设；同样的系统提示与证据拼接口径；同样的高风险标签和任务类别字段；同样的比较对象定义。如果这些边界不一致，就会出现一个非常隐蔽的问题：策略模型在回答一种任务，奖励模型却像在给另一种任务打分。

因此，一个务实建议是：冷启动阶段优先让奖励模型和策略模型共享尽可能多的输入协议，而不是一开始就追求花哨的异构组合。等到团队已经能稳定识别“奖励模型到底错在哪些切片”时，再考虑是否换更轻或更强的奖励架构。先把输入边界统一，常常比先争论模型架构更能降低排障成本。

\subsection{优势函数与 value model：PPO 不是只看“谁分高”}
旧稿在 PPO 细节里花了不少篇幅讲优势函数，这里保留一个工程读者真正需要的版本。PPO 不是简单地看到高奖励回答就无脑加大概率，而是会问一句：\textbf{这次回答比“按当前水平本来该有的结果”好多少？}

把这个直觉写成最小公式，可以记成：
\[
\hat{A}_t \approx R_t - V(s_t),
\]
其中 \(R_t\) 可以理解成这次 rollout 最后拿到的回报，\(V(s_t)\) 则是价值模型对“在当前状态下大概能拿多少分”的估计。于是优势函数更像一个“超出预期多少”的信号，而不是“绝对分有多高”的信号。

这一步在工程上非常重要。因为语言任务里的奖励本来就噪声很大，如果只看绝对奖励，策略模型很容易被偶然高分、格式偏置或奖励模型误判带着跑；有了 value baseline，更新会更像是在问“哪些行为稳定地比当前常态更好”，训练方差会小得多。也正因为如此，旧稿里才会反复提醒：\textbf{value model 不是装饰件，它决定 PPO 更新会不会抖得太厉害。}

\subsection{Actor--Critic 的直觉：一个负责出招，一个负责估价}
很多读者看到 PPO 时会觉得 value model 很抽象，其实可以把它和 policy model 分成两个非常朴素的角色。\textbf{actor} 负责“出这一步回答”，\textbf{critic} 负责“估这一步大概值不值得”。前者像执笔写答案的人，后者像一个不断给出“这条路大概会拿多少分”的内部估价器。

\begin{table}[H]
\centering
\small
\begin{tabularx}{\textwidth}{>{\raggedright\arraybackslash}p{2.4cm}>{\raggedright\arraybackslash}p{3.2cm}>{\raggedright\arraybackslash}X}
\toprule
\textbf{角色} & \textbf{在 RLHF/PPO 里主要做什么} & \textbf{如果这一侧失真会怎样} \\
\midrule
Actor / Policy & 生成候选回答，并根据奖励方向更新分布 & 行为直接漂移，可能更会迎合高分模式 \\
Critic / Value & 估计当前状态下的预期回报，帮助算优势 & 更新忽左忽右，训练方差增大，策略更容易抖动 \\
\bottomrule
\end{tabularx}
\caption{Actor--Critic 的工程直觉是“一个负责出招，一个负责估价”}
\label{tab:actor-critic-intuition}
\end{table}

这层直觉很重要，因为它解释了为什么 PPO 不是“拿到奖励就往上推”那么简单。没有 critic，策略模型只能盯着嘈杂的最终奖励做反应；有了 critic，它才更像在比较“这次比常态好多少”。因此，当 RLHF 训练日志里奖励曲线抖动很大时，团队不该只怀疑 policy learning rate，也要反过来看 value 估计是不是已经开始失真。

\subsection{PPO 为什么常用于 RLHF}
PPO（Proximal Policy Optimization）的优势在于更新相对稳定，不会让策略一步跳得过远。它常见的裁剪目标可以写成：
\[
\mathcal{L}_{\mathrm{PPO}} = \mathbb{E}\left[\min\left(r_t(\theta)\hat{A}_t,\ \mathrm{clip}(r_t(\theta),1-\epsilon,1+\epsilon)\hat{A}_t\right)\right],
\]
其中 \(r_t(\theta)\) 是新旧策略概率比，\(\hat{A}_t\) 是优势函数。直观理解就是：\textbf{奖励高的行为可以被加强，但加强幅度必须被限制在一个稳定区间内。}

工程上更重要的不是公式本身，而是它表达的约束思想：对齐训练不是让模型“越偏好越好”，而是要在可控范围内逐步挪动策略。PPO 之所以常见，正是因为 RLHF 的目标很容易诱导模型朝某些表面模式冲得过猛，例如更长、更自信、更多免责声明、更多看似安全但实际没帮助的话。PPO 的价值就在于给这种“冲得太快”的风险加上刹车。

对教材读者而言，可以把 PPO 理解成一种保守更新原则：如果奖励模型告诉你某类回答更好，策略模型可以朝那个方向移动，但不能一下子把原有分布扔掉。这样做的目的不是追求最优雅的理论，而是降低训练不稳定、语言能力退化和奖励黑客被迅速放大的概率。

\subsection{Rollout 与采样纪律：RLHF 里的“训练数据”是模型自己生成的}
PPO 进入 RLHF 后，一个很大的变化是：训练样本不再只是静态数据集里的固定答案，而是策略模型在当前阶段\textbf{自己采样出来的回答}。这意味着 rollout 阶段本身就会决定你后面能学到什么。

如果采样太贪婪，模型总是给出很相似的回答，偏好优化看到的只是一小片行为空间；如果采样太发散，奖励模型又可能不断给极端样本打分，导致训练噪声暴涨。因此 rollout 不是“随便生成一些回答”，而是要守住两条纪律：
\begin{itemize}
  \item \textbf{既要有足够多样性，又不能完全失控}：否则要么学不到新行为，要么全是噪声。
  \item \textbf{采样口径要稳定}：温度、长度上限、停止条件和提示模板一旦频繁变化，奖励信号就很难横向比较。
\end{itemize}

这也是为什么很多 RLHF 工程实现会把 rollout 配置单独版本化。因为一旦 rollout 分布变了，后面看到的奖励变化可能不再来自“模型学得更好了”，而是来自“你让它生成了另一类样本”。

\subsection{rollout 预算不是算力细节，而是训练节奏本身}
很多团队在第一次做 RLHF 时，会把 rollout 预算理解成“算力同学去解决的资源问题”。但从训练角度看，rollout 预算其实直接决定你一轮更新看到的行为覆盖面、奖励分布稳定性和样本新鲜度。也就是说，它不是外围成本项，而是 PPO 是否能学得稳的核心节奏控制器。

一个很粗但很有用的估算式是：
\[
\text{总 rollout token} \approx \text{提示数} \times \text{每题采样条数} \times \text{平均回答长度}.
\]
若再把奖励模型、参考模型和价值模型的前向都算上，真实训练成本还会继续放大。因此，团队不该只问“这轮能不能跑完”，更该问“这轮 rollout 是否覆盖了值得学的行为差异”。

\begin{table}[H]
\centering
\small
\begin{tabularx}{\textwidth}{>{\raggedright\arraybackslash}p{2.9cm}>{\raggedright\arraybackslash}p{3.1cm}>{\raggedright\arraybackslash}X}
\toprule
\textbf{预算维度} & \textbf{主要影响什么} & \textbf{如果配得不合理会怎样} \\
\midrule
每题采样条数 & 同一提示下可比较的候选多样性 & 太少看不到偏好差异，太多又把成本全耗在少量题上 \\
平均回答长度 & 单条轨迹成本、KL 规模、奖励波动 & 太长时训练更贵，也更容易学到“冗长换高分”的坏模式 \\
每轮提示覆盖数 & 分布覆盖面与样本新鲜度 & 太窄会让策略围着少量题打转，泛化很差 \\
rollout 刷新频率 & on-policy 程度与训练可信度 & 刷新太慢时，当前策略继续吃旧样本，更新方向会越来越虚 \\
\bottomrule
\end{tabularx}
\caption{rollout 预算直接塑造 RLHF 的学习节奏，而不仅是显卡账单}
\label{tab:rollout-budget}
\end{table}

这张表真正想告诉读者的是：如果 rollout 设计得不好，后面很多症状都会被误诊。奖励抖动、KL 飘高、长度膨胀，看起来像 PPO 超参数问题，实际上可能只是因为这轮采样题太少、回答太长、样本太旧。因此，rollout 预算应该和学习率、KL 系数一样，被放进训练设计文档里显式管理。

\subsection{探索不能脱离安全边界：rollout 配置本身也是治理点}
RLHF 里的探索常常被说成“给模型一些自由，让它试出更好的回答方式”。这句话只说对了一半。对企业场景来说，探索必须和安全边界一起设计。否则模型探索出来的，不一定是更好的业务行为，也可能是更会碰红线的写法。

因此，rollout 配置除了控制温度、长度和停止条件，还应显式保留几类安全护栏：高风险提示模板不要随意换；红线拒答场景要始终保留一定比例的监控样本；对权限、财务、法务这类问题，采样出来的回答要优先进入人工抽检，而不是直接拿平均奖励判断“模型进步了”。换句话说，\textbf{探索是为了扩大可选行为，不是为了取消业务边界。}

这一点也解释了为什么 RLHF 训练和普通离线微调不同。离线微调更多是在固定数据上吸收模式，RLHF 则是在边生成边试探。只要有“试探”，就必须同时有“护栏”。如果团队只谈探索、不谈护栏，那么 rollout 阶段产生的数据很可能已经把后面的策略更新带偏了。

\subsection{PPO 的稳定性不只靠 clip：还靠样本新鲜度、epoch 数与早停}
旧稿在 PPO 工程部分还反复提到三个常被低估的稳定器。第一是\textbf{样本新鲜度}。PPO 本质上偏向 on-policy，如果 rollout 样本已经来自一版很旧的策略，再拿来给当前策略反复优化，更新方向就会越来越不可信。第二是\textbf{每批样本反复训练的 epoch 数}。epoch 太少，奖励信号吃不进去；epoch 太多，又容易把一小批 rollout 样本记死，导致过拟合当前奖励分布。第三是\textbf{基于 KL 或奖励异常的早停}。当策略突然偏离太快时，硬把这一轮训练跑完，往往只会把漂移放大。

可以把这一段理解成一个简单判断：
\begin{itemize}
  \item rollout 更新得太慢，等于拿“旧世界”的样本教“新世界”的策略；
  \item 单批样本训得太狠，等于把短期偏差当成长期规律；
  \item 没有早停，等于看见车在打滑还继续踩油门。
\end{itemize}

这也是为什么很多 RLHF 训练日志里，除了奖励和任务指标，还会单独盯 KL、回答长度、拒答率、奖励分布尾部和人工抽检结果。PPO 的 clip 只是在目标函数里加刹车，真正让系统稳住的，是整套 rollout 频率、batch 组织、训练轮数和异常停止纪律一起工作。

\subsection{PPO 失稳时先看哪几块仪表盘}
到了工程现场，团队最常问的不是“PPO 的理论假设是什么”，而是“这一轮看起来不对劲，先查哪里”。教材里需要给读者一套足够朴素的排障顺序，不然 RLHF 很容易沦为只会看总奖励曲线的黑箱训练。

\begin{table}[H]
\centering
\small
\begin{tabularx}{\textwidth}{>{\raggedright\arraybackslash}p{2.9cm}>{\raggedright\arraybackslash}p{3cm}>{\raggedright\arraybackslash}X}
\toprule
\textbf{先看到的症状} & \textbf{更可能是哪类问题} & \textbf{第一批该查的仪表盘} \\
\midrule
奖励快速上涨，但人工抽检变差 & 奖励黑客或 KL 太松 & KL 曲线、回答长度、拒答率、高分样本人工复核 \\
KL 突然飙升 & 策略更新过猛或 rollout 分布突变 & 学习率、采样配置版本、早停阈值、参考模型是否一致 \\
训练很稳但几乎没收益 & KL 太紧或奖励信号太弱 & KL 系数、优势分布、奖励方差、偏好数据区分度 \\
回答越来越模板化 & 奖励偏爱表面特征或采样过保守 & 奖励高分样本形态、temperature、repetition 现象、长度分布 \\
不同批次波动极大 & 样本太旧、batch 太小或高风险样本比例失衡 & rollout 刷新频率、batch 组成、切片分布、value loss 抖动 \\
\bottomrule
\end{tabularx}
\caption{PPO 排障先看仪表盘，而不是先换一套新算法}
\label{tab:ppo-debug-dashboard}
\end{table}

这套仪表盘思维的关键，是把“模型表现不对”重新拆回可观测信号。因为 RLHF 的问题很少只在一个点上发生。它更常见的形态是：偏好数据有一点歪，奖励模型放大一点，rollout 再把噪声扩散一点，最后才由策略模型把问题完整演出来。先建立仪表盘，团队才有可能定位是哪一环先开始失真。

\subsection{KL 约束的意义}
在 RLHF 中常加入 KL 约束，目的是避免策略为了追求奖励而完全偏离原有语言能力。工程上这意味着：我们希望模型“更符合偏好”，但不希望它因此失去可读性、事实性和基础任务能力。

如果把 KL 的作用说得更直白一点，它相当于在提醒训练过程：\textbf{可以更偏向我们喜欢的回答，但不要为了拿高分把原本会的东西也练坏。}

KL 系数之所以重要，是因为它直接控制“往奖励方向走多远”。系数太小，约束太弱，模型可能迅速学会迎合奖励模型的漏洞；系数太大，约束太强，训练几乎推不动策略，看起来做了 RLHF，实际上行为变化非常有限。

\begin{table}[H]
\centering
\small
\begin{tabularx}{\textwidth}{>{\raggedright\arraybackslash}p{2.6cm}>{\raggedright\arraybackslash}X>{\raggedright\arraybackslash}X}
\toprule
\textbf{KL 状态} & \textbf{常见现象} & \textbf{工程后果} \\
\midrule
系数过小 & 奖励快速上升，回答风格突然变长、变硬、变模板化 & 容易奖励黑客，基础能力和真实可用性下降 \\
系数适中 & 奖励逐步提升，回答变化可解释，离线评测与人工抽检一致改善 & 策略稳步靠近业务偏好，回归风险可控 \\
系数过大 & 奖励提升缓慢，回答几乎与 SFT 阶段相同 & 训练成本高，但用户几乎感受不到收益 \\
\bottomrule
\end{tabularx}
\caption{KL 系数过小、适中、过大时的典型后果}
\label{tab:kl-effects}
\end{table}

判断 KL 是否合适，不能只看一个总奖励数值，而要看多个信号是否一致：离线奖励是否上升、人工抽检是否认为更好、任务成功率是否持平或提升、幻觉率是否恶化、回答长度是否异常膨胀、拒答率是否突然升高。若奖励上涨但人工评测变差，往往意味着 KL 太松或者奖励模型本身出了问题。

在企业知识助手里，KL 约束还有一个很现实的作用：保住系统已有的“业务语感”。很多团队花了很久才把格式、术语和知识库引用方式打磨稳定。如果 RLHF 一轮更新就把这些习惯冲散，业务方通常不会认为这是一次“更对齐”的升级，而会认为系统变得更陌生、更难信任。

\begin{figure}[H]
\centering
\begin{tikzpicture}[
  font=\small,
  box/.style={llm box, minimum width=2.8cm, minimum height=1cm},
  arr/.style={llm arr}
]
\node[box] (sft) {SFT 策略模型};
\node[box, right=1cm of sft] (sample) {采样回答};
\node[box, right=1cm of sample] (reward) {奖励模型\\偏好打分};
\node[box, right=1cm of reward] (ppo) {PPO 更新\\含 KL 约束};
\draw[arr] (sft) -- (sample);
\draw[arr] (sample) -- (reward);
\draw[arr] (reward) -- (ppo);
\draw[arr] (ppo) to[bend left=20] node[above]{更新后策略} (sft);
\end{tikzpicture}
\caption{RLHF/PPO 的最小闭环}
\label{fig:rlhf-loop}
\end{figure}

\subsection{奖励黑客与错配风险}
只要奖励模型学得不对，策略模型就可能学会“迎合评分规则”，而不是真正提高可用性。这就是常说的奖励黑客风险。奖励黑客并不神秘，本质上就是优化目标与真实目标不一致时，系统会优先去钻可优化的空子。

例如，奖励模型如果过度偏爱长答案，策略模型就可能在材料不足时也写出冗长、肯定、看似负责但实际风险更高的回答。对企业知识助手来说，这种错配尤其危险，因为它会让系统在最该保守的时候反而更会“装懂”。

常见的奖励模型失败模式至少包括以下几类：
\begin{itemize}
  \item \textbf{表面特征替代真实质量}：把“更长”“更礼貌”“更像模板”误当成“更好”。
  \item \textbf{上下文忽略}：没有真正理解问题风险等级，对高风险和低风险问题使用同一偏好尺度。
  \item \textbf{证据错配}：偏爱看似有引用的回答，却没核查引用是否真的支持结论。
  \item \textbf{分布外脆弱}：在训练样本中表现正常，一到新部门、新制度或新问法就打分漂移。
  \item \textbf{标注员风格污染}：学到的是某批标注员的语气和习惯，而不是组织真正想要的行为。
\end{itemize}

可以把奖励黑客理解为一个三步链条：先是偏好定义不清，接着奖励模型学到表面代理特征，最后策略模型把这些代理特征放大成稳定行为。到了第三步，问题常常会被误诊成“RL 训练不稳定”，但更深层的原因其实是前面两步没有定义好奖励。

一个典型企业场景是报销助手。若奖励模型把“给出明确结论”学成高分特征，策略模型就会倾向于直接回答“这张发票可以报销”。但真实偏好可能是：\textbf{先检查制度依据，再列出缺少字段，最后提醒是否需要人工审批。} 如果奖励模型没学到这种结构，策略优化越成功，系统离真实业务流程反而越远。

因此，奖励模型上线前后都应有针对性诊断：
\begin{itemize}
  \item 看高分回答是否异常趋同，例如都变成长回答或统一模板。
  \item 看证据不足场景下，模型是否更会承认不确定，而不是更会强行下结论。
  \item 看不同风险等级的问题上，模型是否保留了必要的边界差异。
  \item 看离线高分样本在人工复核时是否真的“更好”，而不是“更像会得高分的样子”。
\end{itemize}

\subsection{LIMA 与 RAFT：轻量路线提醒我们别急着上完整 RL 闭环}
在很多真实团队里，完整 RLHF 之前往往还会有两类更轻的路线值得尝试。第一类是 LIMA 风格的“高质量小规模 SFT”。它的重点不是让你永远不做偏好优化，而是提醒你：如果当前问题主要是交互样式、回答骨架和任务口径还没训稳，那么先把高质量 SFT 做好，收益往往比直接进入 RL 更大。

第二类是 RAFT 一类“先用奖励排序，再做监督式吸收”的路线。它不像 PPO 那样维护完整 RL 闭环，而更像先让奖励模型对候选回答排序，再把排序结果回灌进更接近监督微调的训练流程。对工程团队来说，这种路线的现实价值在于：\textbf{保留一部分偏好信号的力量，但把训练链路控制在更短、更容易回归的范围内。}

\begin{table}[H]
\centering
\small
\begin{tabularx}{\textwidth}{>{\raggedright\arraybackslash}p{2.6cm}>{\raggedright\arraybackslash}p{3cm}>{\raggedright\arraybackslash}X}
\toprule
\textbf{轻量路线} & \textbf{更适合什么局面} & \textbf{它提醒你的核心判断} \\
\midrule
LIMA 风格高质量 SFT & 基础任务还没训稳、偏好问题尚未独立成层 & 先把交互样式和基础行为训好，再谈昂贵对齐 \\
RAFT 一类奖励排序微调 & 已有奖励信号，但不想承担完整 RL 闭环复杂度 & 偏好优化可以先走短链路，不必一上来就 PPO \\
\bottomrule
\end{tabularx}
\caption{轻量对齐路线的价值，常常在于帮助团队确认“问题到底到了哪一层”}
\label{tab:light-alignment-routes}
\end{table}

这两条路线之所以值得写进教材，不是因为它们一定比 RLHF 更强，而是因为它们逼着团队先回答一个更重要的问题：\textbf{我们现在缺的是偏好优化能力，还是缺一套足够稳的基础行为基线。}

\subsection{离线偏好优化与在线 RL，是两种完全不同的投资}
学到这里，读者很容易把 DPO、RAFT、RRHF 和 RLHF/PPO 都看成“偏好优化的不同口味”。但从工程投入上看，它们其实分属两类完全不同的投资。前者更像\textbf{离线偏好吸收}：你先把偏好样本整理好，再在相对静态的数据上更新模型。后者则是\textbf{在线策略试探}：模型边生成、边打分、边根据当前策略继续往前挪。

\begin{table}[H]
\centering
\small
\begin{tabularx}{\textwidth}{>{\raggedright\arraybackslash}p{3.1cm}>{\raggedright\arraybackslash}p{3.1cm}>{\raggedright\arraybackslash}X}
\toprule
\textbf{路线} & \textbf{更像在解决什么} & \textbf{什么时候更值得优先上} \\
\midrule
高质量 SFT / DPO / RRHF / RAFT & 把已经看清的偏好稳定写进模型 & 偏好样本已较清楚，但训练基础设施与治理能力仍有限 \\
RLHF / PPO & 在偏好空间里继续做动态探索和保守更新 & 偏好定义成熟，且团队能持续监控奖励、KL、灰度效果与回滚风险 \\
\bottomrule
\end{tabularx}
\caption{离线偏好优化和在线 RL 的成本结构与治理要求并不相同}
\label{tab:offline-vs-online-alignment}
\end{table}

这条区分非常重要，因为很多团队并不是“算法上做不到 RLHF”，而是“治理上还接不住 RLHF”。如果你们还没有稳定的评测门禁、坏例回流、奖励校准和版本回滚能力，那么先把偏好通过 DPO 或高质量 SFT 吸收进去，往往是更负责任的做法。只有当团队不仅能训练，还能监控、解释和回滚时，在线 RL 的价值才真正开始显现。

\subsection{什么时候 RRHF、DPO 或 RLAIF 可能比 RLHF 更现实}
RLHF 并不是唯一的对齐路线。在很多团队里，RRHF、DPO（Direct Preference Optimization）或 RLAIF（Reinforcement Learning from AI Feedback）反而更现实。原因不在于它们一定更先进，而在于它们对工程条件的要求不同。

RRHF 可以理解为一种更偏“排序优化”的中间路线。它仍然利用偏好比较，但不把系统完整拆成“奖励模型 + 强化学习闭环”。如果团队已经有比较清楚的胜负样本，想尽快把“哪个回答更好”的排序倾向写进模型，却又不想立刻承担完整 RLHF 的链路复杂度，那么 RRHF 往往是一个比 RLHF 更轻、比纯 SFT 更贴近偏好优化的选择。

DPO 的吸引力在于：它仍然使用偏好数据，但不需要单独训练一个奖励模型再跑完整 RL 闭环，而是直接利用偏好对策略进行优化。对资源有限、训练基础设施还不成熟的团队而言，DPO 往往更容易落地，因为链路更短、超参数更少、排障成本更低。若团队已经有一批质量不错的成对偏好样本，但还没有稳定的 RL 训练平台，DPO 往往是比 RLHF 更现实的第一步。

RLAIF 的现实意义则在于降低人工偏好标注成本。它不是完全不要人，而是先由更强模型或规则系统生成偏好判断，再由人类做抽检、校准和重点场景把关。对于企业场景中大量中低风险样本，这种方法可以显著降低冷启动成本。但它的前提是：用来给反馈的 AI 评审器本身要足够可靠，否则只是把偏差从一个模型搬到另一个模型。

\begin{table}[H]
\centering
\small
\begin{tabularx}{\textwidth}{>{\raggedright\arraybackslash}p{2.4cm}>{\raggedright\arraybackslash}p{3.2cm}>{\raggedright\arraybackslash}p{3.4cm}>{\raggedright\arraybackslash}X}
\toprule
\textbf{路线} & \textbf{更适合的前提} & \textbf{主要优点} & \textbf{主要风险} \\
\midrule
RLHF & 偏好定义清楚，奖励建模和 RL 基础设施较成熟 & 能显式分离奖励建模与策略优化，适合做细粒度控制 & 链路长、成本高，容易在奖励模型处失真 \\
RRHF & 已有稳定的胜负或排序样本，但不想上完整 RL 闭环 & 比 RLHF 更轻，保留偏好优化味道 & 对偏好样本质量仍然高度敏感，细粒度控制力较弱 \\
DPO & 已有成对偏好数据，但不想承担完整 RL 闭环复杂度 & 训练实现相对简单，工程调参与回归成本较低 & 偏好样本质量决定上限，仍可能继承标注偏差 \\
RLAIF & 人工标注昂贵，且可获得较强 AI 评审器辅助 & 扩展速度快，适合大规模预筛与弱监督扩容 & 评审器偏差会被系统性放大，必须做人类校准 \\
\bottomrule
\end{tabularx}
\caption{RLHF、RRHF、DPO 与 RLAIF 的工程取舍}
\label{tab:alignment-alternatives}
\end{table}

一个务实的判断标准是：如果你们当前最大的瓶颈是\textbf{偏好定义不清}，那无论是 RLHF、RRHF、DPO 还是 RLAIF 都救不了你；如果瓶颈是\textbf{人工标注太贵}，RLAIF 可能更现实；如果瓶颈是\textbf{训练基础设施还弱}，DPO 或 RRHF 往往比 RLHF 更先值得尝试。

\subsection{先选路线，再选算法：一条更实用的对齐决策阶梯}
从教材角度看，很多初学者真正困惑的不是“RLHF 和 DPO 数学上有什么区别”，而是“我们团队现在到底该走哪一条路”。比起先背算法名，更实用的是先走一遍决策阶梯：
\begin{enumerate}
  \item 如果主要问题是知识缺失、检索失败或工具不可靠，先修系统，不进入对齐训练；
  \item 如果主要问题是回答格式、任务骨架和基本交互不稳，先补 Prompt 或 SFT；
  \item 如果偏好问题已经明确，但基础设施较弱，先考虑 RRHF 或 DPO 这种短链路方法；
  \item 如果人工标注成为主要瓶颈，但有可用的强评审器，考虑 RLAIF 做预筛与扩容；
  \item 只有当偏好定义清楚、回归评测稳定、奖励建模与训练基础设施都具备时，再进入完整 RLHF/PPO。
\end{enumerate}

这条阶梯的价值在于把“算法选择”重新放回“问题分层”之后。工程上最昂贵的错误，不是选了某个次优算法，而是在还没把问题层级看清时，就贸然上了最复杂的算法。

\subsection{什么时候不该上 RLHF}
对很多工程团队来说，更重要的问题不是“如何做 RLHF”，而是“现在到底该不该做”。以下几类场景通常都不应优先进入 RLHF：
\begin{itemize}
  \item 问题主要来自知识缺失、文档过期或检索失败。
  \item 连基础 SFT、提示模板和评测集都还不稳定。
  \item 没有可复用的偏好标注数据，只能凭印象判断“回答更好”。
  \item 团队无法做回归评测，不知道行为变化是提升还是漂移。
  \item 业务方并未对“更好”形成共识，不同人偏好彼此冲突。
\end{itemize}

\begin{table}[H]
\centering
\small
\begin{tabularx}{\textwidth}{>{\raggedright\arraybackslash}p{3.1cm}>{\raggedright\arraybackslash}p{2.8cm}>{\raggedright\arraybackslash}p{3.1cm}>{\raggedright\arraybackslash}X}
\toprule
\textbf{观察到的症状} & \textbf{更像哪一层问题} & \textbf{更优先的动作} & \textbf{是否应上 RLHF} \\
\midrule
答案依赖最新制度，但模型答错 & 知识接入问题 & 先修 RAG、知识库更新和引用链路 & 否 \\
回答格式不稳、风格混乱 & 任务行为问题 & 先改提示词、模板或补 SFT & 否 \\
工具调用经常失败，导致动作建议错误 & 系统编排问题 & 先修工具可靠性和异常处理 & 否 \\
多种回答都能答对，但保守性、引用习惯不稳定 & 对齐问题 & 收集偏好数据，再考虑 DPO 或 RLHF & 是，候选 \\
没有稳定评测集，只能凭主观印象判断升级效果 & 基础工程问题 & 先建评测体系和回归流程 & 否 \\
高风险问题上业务方都说不清什么叫更好 & 需求定义问题 & 先统一偏好标准和升级流程 & 否 \\
\bottomrule
\end{tabularx}
\caption{“现在不该上 RLHF”的诊断表}
\label{tab:when-not-rlhf}
\end{table}

这张表的重点不是否定 RLHF，而是强调顺序。RLHF 是处理偏好问题的后手，不是处理所有问题的第一响应动作。越是高成本、长链路的优化方法，越应该在问题已经被正确分层以后再使用。

\subsection{强化学习在 NLP 中的角色}
在大模型之外，强化学习在 NLP 中还常见于：
\begin{itemize}
  \item 对话策略优化。
  \item 长序列决策与工具调用。
  \item 多步推理中的行为规划。
  \item 从执行结果或环境反馈中学习。
\end{itemize}

这提醒我们：RLHF 不是“只有大模型才需要的神秘技术”，它只是把强化学习的一套思想，带到了语言行为优化里。对于 NLP 来说，RL 的价值不只在于“把回答说得更讨喜”，还在于处理\textbf{有延迟反馈、需要多步决策、目标不能完全写成标准答案}的任务。

例如，一个系统如果要连续完成“检索制度、比较条款、调用审批工具、生成最终回复”四步动作，那么它的质量不再只由最后一句自然语言决定，而由整条行为链决定。此时 RL 思想的重要性就会提高，因为你优化的不再只是文本本身，而是整个决策过程。

\section{主案例推进}
延续上一章的评测切片，企业知识助手到了这一章，最值得做对齐的通常不是“产品 FAQ 会不会答”，而是那些\textbf{多个答案都可能不算错、但组织明显更偏好其中一种写法}的高风险场景。也正因为如此，本章主案例继续围绕制度条款查证、权限流程说明和证据不足拒答这几类企业问题来推进。

对于企业知识助手，很多关键行为并不是“事实对错”问题，而是“输出偏好”问题。例如：
\begin{itemize}
  \item 材料不足时，宁愿保守也不要猜测。
  \item 优先给出依据，再给建议动作。
  \item 对高风险流程问题，优先引导人工确认。
  \item 在格式上保持统一，便于工单系统或知识系统接入。
\end{itemize}

这些要求可以先通过提示工程和 SFT 处理；只有当系统规模扩大、偏好边界复杂、且已有足够反馈数据时，RLHF 才值得进入路线图。换句话说，\textbf{对齐训练通常是中后期优化手段，而不是第一步。}

如果把企业知识助手的偏好训练路线再写得具体一点，可以是：
\begin{enumerate}
  \item 先用 SFT 让模型学会固定输出格式和基础任务行为。
  \item 再从人工审核记录、失败会话和高风险场景中提取偏好样本。
  \item 按“证据优先、保守性、权限边界、可执行性”拆分标注规则。
  \item 训练奖励模型，让系统学会区分“更保守、更有依据、更可执行”的回答。
  \item 用 PPO 在 KL 约束下做策略更新，或在基础设施较弱时先尝试 DPO。
  \item 最后用评测集和人工抽检确认：模型真的更稳了，而不是更会迎合奖励。
\end{enumerate}

下面用一个更完整的业务例子串起来。假设用户问：“供应商未盖章的合同扫描件能否直接走付款流程？”现有知识库写明“付款前需完成合同审核”和“例外付款需经法务与财务负责人共同审批”，但没有直接说“未盖章扫描件是否可直接付款”。在这种情况下：
\begin{itemize}
  \item 一个仅追求任务完成度的模型，可能直接回答“不能付款”或“可以走例外流程”。
  \item 一个更符合偏好的模型，应先说明现有依据只覆盖合同审核与例外审批要求，尚不足以直接判断“未盖章扫描件”的合法性。
  \item 它随后应给出更稳妥动作，例如“请补查合同管理制度的签章条款，或升级给法务确认；在确认前不建议直接发起付款”。
\end{itemize}

这个例子说明，企业偏好往往不奖励“说得最像专家”，而奖励“在证据边界内做出最稳妥、最可执行的回应”。如果团队用这类场景持续积累偏好对，奖励模型就有机会学到真正重要的组织偏好，而不是学到表面语气。

在这个案例里，真正的难点并不是写出 PPO 代码，而是把“什么叫更好的企业回答”定义清楚。定义不清，对齐训练就会成为高成本的放大器，把原本就模糊的目标放得更模糊。

\subsection{对齐上线也要先离线后灰度，而且仍需规则系统托底}
即使一轮对齐训练在线下评测里看起来收益明显，也不应直接全量替换线上版本。原因和上一章一致：离线偏好提升解决的是“候选版本值不值得发”，真正发布后还要面对真实问法、真实权限状态、真实并发负载和真实组织流程。因此，对齐版本的上线顺序同样应该是：先离线回归，再小流量灰度，再逐步放量。

但这里还有一个比灰度更重要的边界：\textbf{对齐训练可以改变模型偏好，不能代替规则、权限和审计系统。} 比如 RLHF 可以让模型更倾向于说“请升级给法务确认”，却不能替代真正的审批流；它可以让模型更少越权建议，却不能替代权限校验；它可以让回答更重视证据，却不能替代审计留痕。

\begin{table}[H]
\centering
\small
\begin{tabularx}{\textwidth}{>{\raggedright\arraybackslash}p{2.8cm}>{\raggedright\arraybackslash}p{3.1cm}>{\raggedright\arraybackslash}X}
\toprule
\textbf{能力或责任} & \textbf{对齐训练能改善什么} & \textbf{仍应由什么机制硬性托底} \\
\midrule
回答风格与保守性 & 让模型更愿意承认不确定、先给依据再给建议 & 业务规则、升级流程、人工复核制度 \\
证据优先与引用习惯 & 让模型更倾向于引用材料而非空口结论 & 引用核查、知识库版本治理、日志留存 \\
权限边界意识 & 让模型更少直接越权给结论 & 身份鉴权、权限校验、审批系统 \\
上线风险控制 & 让候选版本更符合偏好目标 & 灰度发布、回滚开关、审计与监控告警 \\
\bottomrule
\end{tabularx}
\caption{对齐训练改善的是偏好，硬边界仍需规则与治理系统托底}
\label{tab:alignment-vs-governance}
\end{table}

教材里把这件事讲清楚很重要，因为很多团队第一次接触 RLHF 时，容易产生一种错觉：只要偏好训得够好，安全和合规问题就会一起消失。现实正相反。越是高风险系统，越要把“模型偏好优化”和“规则系统托底”同时做好，而不是让其中一方替代另一方。

\section{常见坑与工程取舍}
\begin{itemize}
  \item \textbf{把 RLHF 当成通用升级按钮}：没有稳定偏好数据和评测体系时，引入 RLHF 风险很高。
  \item \textbf{奖励模型目标模糊}：奖励定义不清，最终会把模型推向不可解释的行为。
  \item \textbf{忽略 KL 约束}：模型可能为了局部高奖励而丢失基础语言能力。
  \item \textbf{在错误层面使用 RL}：如果问题本质是检索失败、知识过期或工具异常，强化学习无法替代系统修复。
  \item \textbf{把离线高分误当线上收益}：离线奖励上升并不自动等于用户体验提升。
\end{itemize}

再补几个更贴近工程现场的风险：
\begin{itemize}
  \item \textbf{偏好标注标准前后不一}：奖励模型学到的是标注员习惯，而不是团队共识。
  \item \textbf{只看离线奖励，不看真实回归样本}：模型可能更会拿分，但没更可用。
  \item \textbf{奖励过度偏爱某种表面特征}：例如更长、更坚定、更模板化的答案。
  \item \textbf{把对齐训练和安全治理混为一谈}：RLHF 能改善偏好，但不能代替权限、审计和规则系统。
  \item \textbf{忽略替代方案}：很多时候先用 DPO、规则校验、审核分流和评测建设，成本收益比比直接上 RLHF 更高。
\end{itemize}

工程取舍的关键是承认：RLHF 不是一条必经路线，而是一种在特定条件下成立的投资。若你已经有稳定偏好样本、明确奖励目标和成熟回归框架，它是值得考虑的；若这些条件还没有齐，优先把问题分层、把评测建好、把数据定义清楚，通常收益更高。

\section{本章练习}
\begin{exercise}[KL 的作用]
为什么 RLHF 中常常需要加入 KL 约束，而不是只追求奖励越高越好？如果 KL 系数过大或过小，会分别带来什么问题？
\end{exercise}

\begin{solution}
如果只追求奖励，模型可能会为了局部高分而{\heiti 偏离原有语言能力、事实性与任务稳定性}。KL 约束的作用是限制策略更新幅度，让模型在“更符合偏好”和“保持原有可读性、稳定性”之间维持平衡。KL 系数过小，策略容易迅速迎合奖励模型漏洞，导致奖励黑客、长度膨胀或过度拒答；KL 系数过大，则更新几乎推不动，训练成本已经付出，但用户几乎感受不到行为改善。
\end{solution}

\begin{exercise}[是否该上 RLHF]
一个企业知识助手目前存在三类问题：一是知识库更新滞后，二是回答格式不统一，三是当多个答案都可能成立时模型有时不够保守。请判断哪些问题适合优先用 RLHF 处理，哪些不适合，并说明原因。
\end{exercise}

\begin{solution}
前两类问题不适合优先用 RLHF。知识库更新滞后属于知识接入问题，应先修 RAG 与文档更新链路；回答格式不统一属于任务行为问题，应先通过提示模板或 SFT 统一格式。只有第三类“多个答案都可能成立，但模型保守性不稳定”的问题才像真正的对齐问题，此时才值得收集偏好数据，并在具备评测体系后考虑 DPO 或 RLHF。
\end{solution}

\section{延伸阅读}
\begin{itemize}
  \item \textbf{Ouyang et al., \emph{InstructGPT}}：理解 RLHF 在指令模型中的早期经典流程。
  \item \textbf{Schulman et al., \emph{PPO}}：理解 PPO 为什么被广泛用于稳定策略更新。
  \item \textbf{Rafailov et al., \emph{DPO}}：适合理解为什么可以用偏好数据直接优化策略，而不必总是走完整 RLHF 闭环。
  \item \textbf{RLAIF 与 AI feedback 实践资料}：适合理解在人工标注昂贵时，如何用 AI 反馈做弱监督扩容，同时避免评审器偏差被放大。
  \item \textbf{偏好标注与奖励建模实践资料}：适合理解“奖励定义清楚”为什么比“算法更复杂”更重要。
\end{itemize}

\section{小结}
本章建立了对齐训练的工程坐标：为什么需要对齐，RL 的最小概念是什么，RLHF 如何利用偏好数据与奖励模型推进策略更新，PPO 与 KL 约束为何重要，以及强化学习在 NLP 中到底扮演什么角色。更关键的是，本章强调了一个务实判断：真正成熟的团队不会盲目上 RLHF，而是先明确偏好目标、奖励来源、失败模式和回归评测，再决定是用 RLHF、DPO、RLAIF，还是先把更基础的问题修好。

\section{下一章过渡}
不论是否进入 RLHF，只要系统准备上线，就必须面对另一个更现实的问题：延迟、显存和吞吐。下一章将进入推理与性能优化，讨论如何让企业知识助手既可用又可部署。
